
























































































































































\section{Introduction}
Let $G = (V,w)$ be a weighted undirected $n$-vertex graph where $w$ is an $\binom{n}{2}$-dimensional nonnegative real vector assigning a (possibly zero) weight to each edge slot.  For a
nontrivial subset $\emptyset \ne X \subsetneq V$, let $\Delta(X)$ be the set of edges of $G$ with one endpoint in $X$ and one endpoint in $\bar{X} = V \setminus X$.  A \emph{cut} $S$ in $G$ is a subset of edges of the form $\Delta(X)$ for
a nontrivial set $X$.  The sets $X$ and $\bar{X}$ are called the \emph{shores} of the cut.  For a cut $S$, its \emph{weight} is the sum of the weights of the edges in $S$, denoted $w(S)$.
The minimum cut problem is to find the minimum of $w(S)$ over all cuts $S$.  The study of algorithms for the minimum cut
problem in theoretical computer science goes back at least to the 1960's and has given rise to a vast and beautiful literature.
Minimum cut is also a problem of great practical importance with applications
to, for example, clustering algorithms and evaluating network reliability. Randomized algorithms can solve the minimum cut problem in nearly linear time: in 1996 Karger gave an
algorithm with running time $O(m \log^3(n))$ to compute the minimum cut of a weighted graph with $m$ edges \cite{Karger00}.  This was the best known bound until very recently when two independent
works improved on it.  Gawrychowski, Mozes, and Weimann \cite{GMW20} gave a randomized algorithm with running time $O(m\log^2(n))$ \cite{GMW20} and Mukhopadhyay and Nanongkai \cite{MN20} gave a randomized algorithm with time complexity $O(m\frac{\log^2(n)}{\log\log n}+n\log^6(n))$.   Gawrychowski, Mozes, and Weimann \cite{GMW20b} later improved the running time of the Mukhopadhyay and Nanongkai algorithm to $O(m\frac{\log^2(n)}{\log\log n}+n\log^{3+\varepsilon}(n))$.

For \emph{simple} graphs $G$, randomized algorithms are known with running times $O(m \log(n))$ and $O(m + n \log^3(n))$ \cite{GNT20}.  For simple graphs even nearly linear
time \emph{deterministic} algorithms are known.  Kawarabayashi and Thorup gave an $O(m \log^{12}(n))$ time algorithm \cite{KT19}, which was subsequently improved to
$O(m (\log(n) \log \log n)^2)$ by Henzinger, Rao, and Wang \cite{HRW20}.


Our work spans two aspects of the study of the minimum cut problem.  The first is to query complexity lower bounds on minimum cut.
A natural model in which to study the query
complexity of minimum cut is for algorithms allowed to make \emph{cut queries}.  A cut query algorithm can query any subset $\emptyset \ne X \subsetneq V$ and receives the
answer $w(\Delta(X))$.  One motivation to study cut query algorithms comes from submodular function minimization.  The cut function $f(X) = w(\Delta(X))$ is a submodular
function, and finding the minimum cut value is equivalent to finding the minimum value of $f$ over all nontrivial sets $X$.  The problem of minimizing a submodular function is often studied with respect to an evaluation
oracle, which in the case of the cut function is exactly a cut query.

Harvey \cite{Harvey08} observed that results on the deterministic communication complexity of deciding graph connectivity \cite{HMT88} imply that any deterministic cut query algorithm to compute minimum cut,
or even to decide if the graph is connected or not, must make at least $c n$ cut queries, for a constant $c < 1$.  Analogous results on the randomized communication complexity of connectivity \cite{BFS86}
imply an $\Omega(n/\log(n))$ lower bound on the number of cut queries needed by a randomized algorithm to compute minimum cut (or even connectivity).

On the algorithms side, Rubinstein, Shramm, and Weinberg \cite{RSW18} gave a randomized algorithm computing the minimum cut of a simple graph with $O(n \log(n)^3)$ many
cut queries\footnote{\cite{RSW18} state the bound as $\tilde O(n)$ but we estimate their sparsifier based algorithm to make $O(n \log(n)^3)$ many cut queries.}.
Recently, Mukhopadhyay and Nanongkai \cite{MN20} used a different approach based on Karger's 2-respecting tree algorithm
\cite{Karger00} to also give a randomized $\tilde O(n)$ cut query algorithm to compute minimum cut in a general undirected weighted graph.

For deterministic cut query algorithms, there remains a large gap between the best upper and lower bounds.  We are not aware of any deterministic algorithm for minimum cut
better than learning the entire graph, which can take $\Omega(n^2/\log(n))$ cut queries in the worst case.  On the lower bound side, Graur, Pollner, Ramaswamy, and Weinberg \cite{GPRW20} recently introduced
a very interesting lower bound technique called the \emph{cut dimension}, which we now describe.  Let $G = (V,w)$ be a weighted undirected graph with $n$ vertices and $m$ edges, and let $\M(G)$ be
the set of minimum cuts of $G$.  For a cut $S \in \M(G)$,
let $\chi(S) \in \{0,1\}^{m}$ be the characteristic vector of $S$
amongst the $m$ edges of $G$.  Let $\vec{\M}(G) = \{\chi(S): S \in \M(G) \}$.  The cut dimension of $G$, denoted
$\cd(G)$, is the dimension of $\laspan(\vec{\M}(G))$.
It is shown in \cite{GPRW20} that for any $n$-vertex graph $G$, the cut dimension $\cd(G)$ is a lower bound on the deterministic cut query complexity of computing
minimum cut on weighted $n$-vertex graphs.  Moreover, for every $n \ge 2$ they construct an $n$-vertex graph $G$ with cut dimension $3n/2-2$.

Besides showing lower bounds on cut query complexity, the cut dimension is a natural measure of the complexity of mincuts in a graph.  There is a rich
literature on the possible structure of mincuts in a graph.  Perhaps the first result of this kind is the cactus representation of mincuts by \cite{DKL76}.
A cactus for a graph $G$ is a sparse weighted graph $C$ that represents all the mincuts of $G$.  One consequence of the
cactus representation is that the number of possible mincuts in an $n$-vertex weighted graph is at most $\binom{n}{2}$.  This upper bound was later given
an algorithmic proof via Karger's famous contraction algorithm (\cite{Karger93}, Theorem~6.1).  The
$n$-vertex cycle graph has $\binom{n}{2}$ many minimum cuts and shows that this bound can be tight.

While the cycle has $\binom{n}{2}$ many mincuts, these cuts live in an $n$-dimensional space as the $n$-vertex cycle only has $n$ edges.  Is it possible to construct graphs
with many cuts that also have high cut dimension?  We show that this is not possible, and in fact the cut dimension of an $n$-vertex graph is at most $2n-3$.
\begin{theorem}[Main Upper Bound]
\label{thm:cd_upper}
For any weighted undirected graph $G$ on $n\ge 2$ vertices it holds that $\cd(G) \le 2n-3$.
\end{theorem}
Like the cactus representation, this shows another aspect in which the mincuts of a graph are constrained to have a relatively simple structure.  We further show that this bound is
tight by constructing graphs with cut dimension $2n-3$ for every $n \ge 2$.
\begin{theorem}[Main Lower Bound]
\label{thm:cd_lower}
For every $n \ge 2$ there exists an $n$-vertex weighted undirected graph $G$ with $\cd(G) = 2n-3$.
\end{theorem}



In addition to shedding further light on the structure of minimum cuts, this improves the best known lower bound on the deterministic cut query complexity of
the minimum cut problem to $2n-3$.  We additionally show that the cut dimension is even a lower bound on a stronger query model called the linear query model, recently
studied in \cite{ACK20}.  In the linear query model, the algorithm can query any vector $x \in \R^{\binom{n}{2}}$ and receives the answer $\braket{w}{x}$, the
inner product of $w$ and $x$.  Linear queries can be much more powerful than cut queries as one can completely learn an unweighted graph with a single linear query.
By an information theoretic argument learning an unweighted graph can require $\Omega(n^2/\log(n))$ many cut queries since each cut query reveals at most
$O(\log (n))$ bits.

We further introduce a lower bound technique which is a generalization of the cut dimension that we call the $\ell_1$-approximate cut dimension.  This technique
looks not just at mincuts in the graph, but all cuts.  We again look at the span of the dimension of these cuts with
an additional twist.  Suppose the weight of a minimum cut in $G$ is $\lambda$ and cut $S$ has $w(S) = \lambda + \delta$.  Abusing notation we will let $S$
represent both a set of edges and the characteristic vector $S \in \{0,1\}^{\binom{n}{2}}$ of $S$ among all edge slots.
The vector $S$ can be \emph{perturbed} to $S - u$ for
any vector $u \ge 0$ with $\|u\|_{1,w} \le \delta$.  Here $\| u \|_{1,w} = \sum_i |w(i) \cdot u(i)|$ is the $\ell_1$ norm of $u$ weighted by the edge weights of
the graph.  The $\ell_1$-approximate cut dimension of $G$ is then the minimum over all valid perturbations of the dimension of the span of the perturbed cut vectors.

The minimization over all perturbations makes the $\ell_1$-approximate cut dimension a difficult quantity to lower bound.  We are able to show, however, that the $\ell_1$-approximate cut
dimension can be strictly larger than the cut dimension.  For every $k \in \N$ and $n = 3k+1$, we construct an unweighted $n$-vertex graph $G$ whose $\ell_1$-approximate cut dimension is $2n-2$.
This has the following application.
\begin{theorem}
\label{thm:2n-2}
Any deterministic linear query algorithm that correctly computes the minimum cut of all $n$-vertex weighted undirected graphs must make
at least $2n-2$ queries in the worst case.
\end{theorem}

Computing the minimum cut of a graph with cut queries is a special case of finding the nontrivial minimum of a symmetric submodular function $f:2^V \rightarrow \R$ with evaluation queries.  That is,
to find $\min_{X: \emptyset \ne X \subsetneq V} f(S)$ for a submodular $f$ that satisfies $f(X) = f(V \setminus X)$ for all $X \subseteq V$.  As linear queries are more powerful than cut queries, \cref{thm:2n-2} also
implies a $2n-2$ evaluation query lower bound for a deterministic algorithm finding a nontrivial minimum of a symmetric submodular function, which is currently the best known.































































































































\setcounter{theorem}{3}\section{Preliminaries}\label{sec:prelim}

For every natural number $n$, we denote by $[n]$ the set
$\{1,2, \ldots, n\}$.
For a vector $z \in \R^n$ we write $z \geq \mathbf{0}$ if every coordinate of the vector
is at least 0, and similarly we write $z = \mathbf{0}$ if $z$ is the all-zero vector.
We denote the scalar product of two vectors $z,z' \in \R^n$
by $\braket{{z}}{z'}$. For any matrix, denote the rank of $A$ by $\rk(A)$.
We denote the disjoint union of sets $X$ and $Y$ by $X \sqcup Y.$

\subsection{Graphs, cuts, sets}

An undirected {\em weighted graph} on $n$ vertices is a couple
$G = (V,w)$, where $V$ is the set of vertices with $|V| = n$,
the set of edge slots
$V^{(2)}$ is the
set of subsets of $V$ with cardinality 2, and the weight function
$w : V^{(2)} \rightarrow \R $ is non-negative. We refer to the vertex set of $G$ as $V(G)$.
The set of edges of $G$
is defined as $E = \{e \in V^{(2)} : w(e) > 0\}$.
When in a graph $G=(V,w)$ the weight of every edge is 1, we say that the graph is {\em unweighted},
and we refer to it also as $G=(V,E)$; such graph is also called a {\em simple graph}.
For an edge $e = \{u,v\} $, we say that $u$ and $v$ are the endpoints of $e$.
For a subset $X \subseteq V$ of the vertices, we denote by $E(X)$ the set of edges in $E$
which have both endpoints in $X$, and for disjoint subsets $X, Y \subseteq V$, we denote by
$E(X, Y)$ the set of edges with exactly one endpoint in each of the two sets.
We extend the weight function $w$ to any subset $E'$ of the edges by $w(E') = \sum_{e \in E'} w(e)$.
We will deal only with graphs which have at least $2$ vertices.

We fix an ordering
$v_1 < v_2 < \cdots < v_n$ of the vertices which induces also an ordering $\{v_1, v_2\}, \{v_1, v_3\}, \ldots, \{v_{n-1},v_n\}$ of the edge slots as well as an ordering $e_1 < e_2 < \ldots < e_m$ of the $m=|E|$ edges.  We view $w \in \R^{\binom{n}{2}}$ as a vector whose $\ith$ coordinate gives the (possibly zero) weight of the $\ith$ edge slot
according to this ordering,
and we define $\vec{w} \in \{0,1\}^m$ as the restriction of $w$ to the edges.
With some slight abuse of notation, for a set of edges $S \subseteq E$, we use the same symbol $S$ to also denote
the characteristic vector in $\{0,1\}^{\binom{n}{2}}$ of $S$ among all edge slots.
We further need the characteristic vector of $S \subseteq E$ among the $m$ edges $E$, for which we use the
notation $\chi(S) \in \{0,1\}^m$.
For a family ${\cal F}$ of subsets of the edges, we use the notation
$\vec{{\cal F}} = \{ \chi(S) \in \{0,1\}^m : S \in {\cal F}\}$.


For $X \subseteq V$, we denote by $\bar{X}$ the set $V \setminus X$.
A {\em cut} $S$ is a set $E(X, \bar{X})$ for some $\emptyset \neq X \subsetneq V$.
We call
$X$ and $\bar{X}$ the {\em shores} of $S$, and we denote the cut by $\Delta(X)$.
A cut is a {\em star cut} if one of its shores is a singleton,
otherwise it is {\em non-star} cut. If the singleton shore of a star cut $S$ is $\{v\}$, then we say that $S$ is a star cut
at $v$.  The {\em weight} of a cut is the sum of the weights of its edges.
For a cut $S$ we define the {\em graph of the cut} $S$ as the unweighted graph $G(S)= (V', E')$
where $V'$ is the set of vertices in $V$
that are endpoints of at least one edge in $S$, and $E'=S$.
We say that a cut $S$ is
{\em connected} if $G(S)$ is a connected graph.
A cut is a {\em minimum} cut, or mincut, for short, if no other cut has smaller weight.
We denote by $\M(G)$ be the set of minimum cuts of $G$.
The {\em cut dimension} of $G$ is $\cd(G) = \dim(\laspan(\vec{\cal{M}}(G)))$.

Let $V$ be a set of size $n$. Two sets $X,Y \subseteq V$ are said to \emph{overlap} if $X \cap Y \ne \emptyset, \bar{X} \cap Y \ne \emptyset,
X \cap \bar{Y} \ne \emptyset$.
 A family $\Gcal$ of subsets of $V$ is said to be \emph{laminar} if for all $X,Y \in \Gcal$
it holds that $X$ and $Y$ do not overlap. A set family $\Gcal \subseteq 2^{V}$ is said to be \emph{closed under overlaps}
if for every $X,Y \in \Gcal$ that overlap it holds
that $X \cap Y, X \cup Y \in \Gcal$.  A laminar subset $\Lcal \subseteq \Gcal$ is said to be \emph{maximal in} $\Gcal$ if for every $X \in \Gcal - \Lcal$
there is a $Y \in \Lcal$ such that $X,Y$ overlap.  We say a laminar subset $\Lcal$ is maximal if it is maximal in $2^{V}$.


The sets $X,Y \subseteq V$ \emph{cross} if they overlap and additionally $\bar{X} \cap \bar{Y} \ne \emptyset$.
Note that if $X,Y$ cross then so do $X,\bar{Y}$. A set family $\Gcal \subseteq 2^{V}$ is said to be {\em cross-free} if for all $X,Y \in \Gcal$ it holds
that $X$ and $Y$ do not cross. Observe that if $X$ and $Y$ do not cross
then either $Y$ or $\bar{Y}$ is a subset of $X$ or $\bar{X}$.
Let $G = (V,w)$ be a graph with $n$ vertices.
Two cuts $\Delta(X)$ and $\Delta(Y)$ of $G$ are {\em crossing} if $X$ and $Y$ are crossing.
Let $\Fcal = \{ \Delta(X_1), \ldots,  \Delta(X_k) \}$ be a set of cuts of $G$.
We say that $\Fcal$ is {\em cross-free family of cuts} if
$\Gcal = \{X_1, \ldots, X_k\}$ is cross-free.  Note that it does not matter which shore we take to be in $\Gcal$.

There is a close relationship between cross-free families of cuts and laminar sets.
Let $\Fcal = \{ \Delta(X_1), \ldots,  \Delta(X_k) \}$ be a cross-free family of cuts where each $X_i \subseteq V$,
and let $X_i' = X_i$ if $v_1 \not \in X_i$ and $X_i' = \bar{X}_i$ otherwise.
The \emph{beach} of $\Fcal$ is the set $\Gcal = \{X_1', \ldots, X_k'\}$.
For a family of sets $\Gcal \subseteq 2^{V}$ we say that it is \emph{proper} if $\emptyset, V \not \in \Gcal$,
and we say that it is \emph{complement free} if it does not contain $X,Y$ with $Y = \bar{X}$.

\begin{claim}\label{claim:cross-free}
Let $\Fcal$ be a cross-free family of distinct cuts and $\Gcal$ its beach.  Then $\Gcal$ is proper, complement free and laminar.
\end{claim}

\begin{proof}
First, $\Gcal$ does not contain $\emptyset$ or $V$ because these are not shores of cuts.  It is complement free because $\Fcal$ contains
distinct cuts, and its beach contains exactly one representative shore from each cut.  Finally, we show that it is laminar.  Let
$X_1, X_2 \in \Gcal$.  By definition of a beach, neither of these sets contain $v_1$, thus $\bar{X}_1 \cap \bar{X}_2 \ne \emptyset$.  Therefore if
$X_1, X_2$ overlapped they would also cross, in contradiction to $\Fcal$ being a cross-free family of cuts.
\end{proof}

A mincut is {\em crossless} if no other
mincut crosses it. Observe that a star mincut is always crossless.
Also, if a  mincut $\Delta(X)$ is crossless then
for every mincut $\Delta(Y)$, either $Y$ or $\bar{Y}$ is a subset of $X$ or $\bar{X}$.
Crossing mincuts have a nice structural property which was already observed by~\cite{DKL76}.

\begin{claim}
\label{claim:mincuts}
Let $G=(V,w)$ be a weighted graph.
If $\Delta(X), \Delta(Y) \in \M(G)$ cross then
$\Delta(X \cap Y), \Delta(X \cup Y)  \in \M(G)$.
\end{claim}

\begin{proof}
We have $\Delta(X \cap Y) \neq \emptyset $ and $ \Delta(X \cup Y) \neq V$ because $\Delta(X)$ and $\Delta(Y)$ cross. The cut function is submodular therefore we have
\begin{equation*}
\label{eq:mincuts}
w(\Delta(X \cap Y) ) + w(\Delta(X \cup Y) ) \leq w(\Delta(X)) + w(\Delta(Y)).
\end{equation*}
Let $c$ be the weight of a minimum cut in $G$. Then the right hand side of the above inequality
is equal to $2c$, while its left hand side is at least $2c$. Therefore
$w(\Delta(X \cap Y) ) + w(\Delta(X \cup Y) ) = 2c$ from which the statement follows.
\end{proof}

\subsection{Operations on graphs}
We will use several operations on graphs.  The first of these is the \emph{direct union}.
\begin{definition}[direct union]
\label{def:du}
For two graphs $G_0= (V_0, w_0), G_1 = (V_1,w_1)$
with disjoint vertex sets,
and for vertices $v_0 \in V_0$ and $v_1 \in V_1$, the {\em direct union} of $G_0$ and $G_1$ at vertices
$v_0,v_1$ is the fusion of the two by identifying $v_0$ and $v_1$. Formally, the direct union is
$G_0^{v_0} \oplus   G_1^{v_1} = (V,w)$
where $V = ( V_0 \cup V_1 \cup \{v\} ) \setminus \{v_0,v_1\}$, for a new vertex
$v \not\in V_0 \cup V_1 $. The weight function  of $G_0^{v_0} \oplus   G_1^{v_1}$
is defined by
\begin{equation*}
w (\{x,y\}) =
\begin{cases}
w_b (\{x,y\}  & \text{ if } x,y \in V_b \setminus \{v_b\}, b\in \{0,1\}, \\
w_b(\{x,v_b\}) & \text{ if } x\in V_b \setminus \{v_b\}, y=v, b\in \{0,1\}, \\
0 &  \text{ otherwise }.
\end{cases}
\end{equation*}
\end{definition}
The cut dimension of a direct union is a simple function of the cut dimensions of its components.

\begin{claim}
\label{claim:directsum}
Let $G = G_0^{v_0} \oplus  G_1^{v_1}$ be the direct union of $G_0$ and $G_1$ at vertices
$v_0, v_1$.
Let $c_b$ be the weight of a minimum cut in $G_b$, for $b=0,1$.
Then $\cd(G) = \cd(G_0) + \cd(G_1)$ if
$c_0 = c_1$, and $\cd(G) = \cd(G_b)$ if $c_b < c_{1-b}$.
\end{claim}

\begin{proof}
Let $\Delta(X)$ be an arbitrary cut of $G$ where $v \not \in X$.
If $X \not \subseteq   V_b$,
for $b \in \{0,1\}$, then
the weight of the cut $\Delta(X)$ is at least $c_0 + c_1$, and therefore
it is not a minimum cut.  If $X \subseteq   V_b$, for some $b \in \{0,1\}$ then
the weight of $\Delta(X)$ in $G$ is the same as the weight of $\Delta(X)$ in $G_b$.
Therefore if $c_0 = c_1$ then every mincut in $G_0$ and
every mincut of $G_1$ is a mincut of $G$, and these are the only mincuts. Since their supports are disjoint,
we have $\cd(G) = \cd(G_0) + \cd(G_1)$. If $c_b < c_{1-b}$ then only the mincuts of $G_b$ are mincuts of $G$,
and therefore $\cd(G) = \cd(G_b)$.
\end{proof}

The next two operations, which are inverses of each other, give a decomposition of a graph along a cut into two smaller graphs, and a composition of two graphs into a bigger one by unfolding a star cut in
each components. The decomposition operation was essentially defined in~\cite{FF09}.
Let $G= (V,w)$ be a weighted graph and let $Z$ be a cut in $G$ with shores $X_0$ and $X_1 = V \setminus X_0$.
The {\em separation} of $G$ along the cut $Z$, denoted by $\sep(G,Z)$, is the set of two
graphs $\{ G_0= (V_0, w_0), G_1 = (V_1,w_1) \}$, where $V_b = X_b \cup \{v_{1-b}\}$, for
$b=0,1$ with new vertices $v_0, v_1$. The respective weight functions are defined by
$w_b (\{x,y\}) = w (\{x,y\}) $ for any $x,y \in X_b$, and $w_b (\{x,v_{1-b}\}) = \sum_{y \in V_{1-b}} w(\{x,y\})$ for any $x\in X_b$.

Let $G_0= (V_0, w_0), G_1 = (V_1,w_1)$ be two graphs on disjoint vertex sets, and let $v_b \in V_{1-b}$
be arbitrary vertices for $b \in \{0,1\}$. The {\em merge} of $G_0$ and $G_1$ along the vertices $v_1, v_0$,
denoted by $\mer( \{ (G_0, v_1), (G_1, v_0) \} )$,
is the graph $G=(V,w)$, where $V = (V_0 \cup V_1) \setminus \{v_0, v_1\}$. The weight function in $G$
is defined by $w (\{x,y\}) = w_b (\{x,y\})$ if $x,y \in V_b$, for $b \in \{0,1\}$, and
\[w (\{x,y\})  = w_0 (\{x,v_1\})  w_1 (\{v_0,y\}), \text{ if } x \in V_0 \text{ and }y \in V_1.\]
It follows from the definitions $\sep$ is the left inverse of $\mer$ if the star cut
at $v_1$ in $V_0$ and the star cut at $v_0$ in $V_1$ both have weight one,
and $\sep$ is the right inverse of $\mer$ if the weight of the cut $Z$ is one. We formally state
the former property.

\begin{claim}
\label{claim:sepmer}
Let $G_0= (V_0, w_0)$ and $ G_1 = (V_1,w_1)$ have disjoint vertex sets, and let $v_b \in V_{1-b}$ such that
$w_b(\Delta(v_{1-b})) =1$,
for $b=0,1$. Let $Z$ be the cut in $\mer( \{ (G_0, v_1), (G_1, v_0) \} )$ whose shores are
$V_0  \setminus \{v_1\}$ and $V_1 \setminus \{v_0\}$. Then $w(Z)=1$ and
$$
\sep ( \mer( \{ (G_0, v_1), (G_1, v_0) \} ), Z ) = \{ G_0, G_1 \}.
$$
\end{claim}

\subsection{Query models}
\begin{definition}[$\MINCUT_n$]
The input in the $\MINCUT_n$ problem is an $n$-vertex weighted undirected graph $G = (V,w)$.  The required output on $G$ is the weight
of a minimum cut in $G$.
\end{definition}
A deterministic algorithm correctly solves the $\MINCUT_n$ problem if it outputs the correct mincut weight for every $n$-vertex input graph $G$.
We consider algorithms given two models of query access to the input graph $G = (V,w)$, linear queries and cut queries.
A {\em linear query} for $G$ is a vector $x \in \R^{\binom{n}{2}}$,
and the query is answered by $\braket{x}{w}$.  A cut query is a vector $x \in \{0,1\}^{\binom{n}{2}}$
which is the characteristic vector of a cut in the complete $n$-vertex graph.  The answer
to a cut query is again $\braket{x}{w}$.  Clearly any cut query algorithm can be simulated by a linear
query algorithm.

We use $D_\cut(\MINCUT_n)$ to denote the minimum, over all deterministic query algorithms $\Acal$ that correctly solve $\MINCUT_n$,
of the maximum over all $n$-vertex input graphs $G = (V,w)$ of the number of cut queries made by $\Acal$ on $G$.
$D_\lin(\MINCUT_n)$ is defined analogously for linear queries.

Some authors instead define the output of the minimum cut problem to be a cut $S$ that achieves the minimum weight, rather than the weight itself.  Over $n$-vertex weighted
graphs let us denote this problem as $\ARGMINCUT_n$.  For linear and cut queries, an algorithm that finds a minimum cut $S$ can also return the weight of $S$ with one additional query.
Thus $D_{\lin, \cut}(\ARGMINCUT_n) \ge D_{\lin, \cut}(\MINCUT_n) - 1$, and the lower bounds we prove for $\MINCUT_n$ can be applied, minus 1, to $\ARGMINCUT_n$ as well.

%

\section{Lower bounds on the linear query complexity of MINCUT}\label{sec:min-query}
Graur et al.~\cite{GPRW20} introduce the cut dimension as a means to show lower bounds on the
deterministic cut query complexity of computing minimum cut.

\begin{theorem}[\cite{GPRW20}]
\label{thm:graur_lower}
If there is an $n$-vertex {weighted graph $G = (V,w)$} with $\cd(G) = k$ then $D_\cut(\MINCUT_n) \ge k$.
\end{theorem}

We show that this theorem even holds with respect to a stronger computational model where the algorithm is able to
make linear queries.  We also give a generalization of the cut dimension to a quantity which is at least as large,
and can be strictly larger, that we call the $\ell_1$-approximate cut dimension.  We now give an overview of the Graur et al.\  \cite{GPRW20}
argument in the context of linear queries and how we can extend it.

The proof of \cref{thm:graur_lower} is based on an adversary argument.  Suppose a deterministic cut query algorithm makes $k$ linear queries
and consider the execution of the algorithm on a fixed $n$-vertex graph $G = (V,w)$ whose set of minimum cuts is $\M(G)$.
Make a $k$-by-$\binom{n}{2}$ matrix ${A}$ whose rows are the query vectors asked by the algorithm. Suppose we can find a vector
${z} \in \R^{\binom{n}{2}}$ such that
\begin{enumerate}
  \item $w-{z} \ge \mathbf{0}$,
  \item ${A}{z} = \mathbf{0}$,
  \item There is a cut $S \in \M(G)$ such that $\braket{S}{z} > 0$.
\end{enumerate}
The existence of such a vector $z$ means the algorithm cannot correctly compute minimum cut
weight on all weighted $n$-vertex graphs.  The reason is that $G' = (V,w-z)$
is a valid non-negatively weighted graph by (1), has the same answers on all queries asked by the algorithm by (2),
and by (3) has minimum cut weight at most $\braket{S}{w-z} = \braket{S}{w} - \braket{S}{z} < \braket{S}{w}$,
which is strictly less than the minimum cut weight of $G$. As with $k$ queries the algorithm cannot distinguish whether the input is $G$ or $G'$, it cannot correctly output the minimum cut weight for all $n$-vertex weighted graphs.

A weaker condition than~(3) suffices for this argument to work.  Suppose that the minimum cut weight in $G$ is $c^*$.
Then the argument still goes through with the condition
\begin{itemize}
  \item[3'.] There is a cut $S$ such that $\braket{S}{z} > \braket{S}{w}-c^*$.
\end{itemize}
This is because the algorithm cannot distinguish the graph $G$ with minimum cut weight $c^*$ from the graph $G' = (V,w-z)$
which has minimum cut weight at most $\braket{S}{w-z} < c^*$.

In order to understand what kind of bound this argument gives, for fixed $w,A,S$ we define the quantity $\alpha(w,A,S)$
which is given by the following linear program.
\begin{equation*}
\begin{aligned}
\alpha(w,A,S) =\  & \underset{z}{\text{maximize}}
& & \braket{S}{z} \\
& \text{subject to}
& & w-z \ge 0 \\
& & & Az = \zero
\end{aligned}
\end{equation*}
Taking the dual of this program gives
\begin{equation*}
\begin{aligned}
\alpha(w,A,S) =\  & \underset{v}{\text{minimize}}
& & \braket{S - A^Tv}{w} \\
& \text{subject to}
& & S - A^T v \ge 0
\end{aligned}
\end{equation*}

The dual tells us that a vector $z$ having large overlap with $S$ and satisfying items $(1),(2)$ above exists
iff the vector $S$ is \emph{far away} from the rowspace of $A$.  The notion of far away here is a one-sided $\ell_1$ distance
weighted by $w$.  It is one-sided because the condition $S - A^T v \ge 0$ tells us we are looking to approximate $S$ by
vectors in the rowspace of $A$ that are entrywise at most $S$.  As $S - A^T v \ge 0$ and $w \ge 0$ this means
$\braket{S - A^Tv}{w} = \sum_i | w(i) \cdot (S (i)- A^Tv)| = \|S - A^Tv \|_{1,w}$, where
$\|u\|_{1,w}$ is defined to be $\sum_i |u(i)w(i)|$.  Thus the value of the dual can be interpreted as the one-sided $\| \cdot \|_{1,w}$ distance
between $S$ and the rowspace of $A$.

This leads us to define an $\ell_1$ approximate version of the cut dimension.  The notion
we need is given by the following definitions.
%
\begin{definition}[one-sided row-by-row $\ell_1$-approximate rank]
Let $Y \in \R^{M \times N}$ be a matrix, $w \in \R^N$ a weight vector and
$c \in \R^M$ a cost vector.  We define the $(w,c)$ one-sided row-by-row $\ell_1$-approximate rank of $Y$ to be
the minimum rank of a matrix $\tilde Y$ such that $\tilde Y \le Y$ and $\| Y(i,:) - \tilde Y(i,:) \|_{1,w} \le c(i)$,
for all $1 \leq i \leq M$.
\end{definition}

\begin{definition}[$\ell_1$-approximate cut dimension]
Let {$G = (V,w)$} be an $n$-vertex weighted undirected graph with minimum cut weight $c^*$.
Let $M$ be $(2^{n-1}-1)$-by-$\binom{n}{2}$ matrix whose rows are $S \in  \{0,1\}^{\binom{n}{2}}$ for all cuts $S$ of $G$.
Let $c = Mw - c^*\1$, where $\1$ is the all one vector.  Then the $\ell_1$-approximate cut dimension of $G$, denoted $\widetilde{\cd}(G)$, is
the $(w,c)$ one-sided row-by-row $\ell_1$-approximate rank of $M$.
\end{definition}

\begin{theorem}
\label{thm:approx_cut}
If there is an $n$-vertex graph {weighted graph $G = (V,w)$} with $\widetilde{\cd}(G) = k$
then $D_\lin(\MINCUT_n) \ge k$.
\end{theorem}

\begin{proof}
Let $G = (V,w)$ be a graph with $\widetilde{\cd}(G) = k$ and let $c^*$ be the minimum cut weight of $G$.  Suppose for
contradiction there is a deterministic $k-1$ linear query algorithm that correctly computes the minimum cut of any
$n$-vertex graph.  Run this algorithm answering queries according to $G$ and
package the queries into a $(k-1)$-by-$\binom{n}{2}$ matrix $A$.

As the algorithm is correct, for every cut $S$ of $G$ it must be the case that $\alpha(w,A,S) \le \braket{S}{w} -c^*$.  If not, the graph $G' = (V,w-z)$, where $z$ is an optimal
solution to the primal of $\alpha(w,A,S)$, has minimum cut weight strictly smaller than $c^*$, yet $G'$ cannot be distinguished from $G$ by the algorithm.
Thus by the dual formulation of $\alpha(w,A,S)$, this means that for every cut $S$ of $G$ there is a vector $\tilde{S}= A^T v$ in the rowspace of $A$ such
that $\tilde S \le S$ and $\|S-\tilde{S}\|_{1,w} \le \braket{S}{w} -c^*$.  The matrix $\tilde M$ whose rows are $\tilde S$ for all
cuts $S$ therefore witnesses that $\widetilde{\cd}(G) \le \rk(A) \le k-1$, a contradiction.
\end{proof}

\begin{lemma}
\label{lem:cdless}
For any weighted graph $G = (V,w)$ we have $\cd(G) \le \widetilde{\cd}(G)$.
\end{lemma}

\begin{proof}
Suppose that $G = (V,w)$ has minimum cut weight $c^*$, and let $\M(G)$ be the set of minimum cuts of $G$.
Let $M$ be the $(2^{n-1}-1)$-by-$\binom{n}{2}$ matrix whose rows are $S \in \{0,1\}^{\binom{n}{2}}$ for all cuts $S$ of $G$
and let $c = Mw-c^*$.

Let $Y$ be the submatrix of $M$ where rows are restricted to cuts in $\M(G)$ and columns are restricted
to the edge slots $e$ where $w(e) > 0$.  Thus the rows of $Y$ are exactly the vectors $\chi(S)$ for $S \in \M(G)$.
and the rank of $Y$ is $\cd(G)$.  Any matrix $\tilde M$ which satisfies $\tilde M \le M$ and $\|M(i,:)-\tilde{M}(i,:)\|_{1,w} \le c(i)$
for all $i$ must contain $Y$ as a submatrix, as $c(i)=0$ for rows $i$ that correspond to minimum cuts and $w$ is positive on the
edge slots labeling the columns of $Y$.  Thus $\rk(\tilde M) \ge \rk(Y)$ for any $(w,c)$ one-sided row-by-row $\ell_1$ approximation
$\tilde M$ of $M$, giving the lemma.
\end{proof}

In \cref{sec:approx_construction} we will see that $\widetilde{\cd}(G)$ can be strictly larger than $\cd(G)$.  From \cref{thm:approx_cut} and
\cref{lem:cdless} we obtain the following corollary.
\begin{corollary}
\label{cor:cd_linear}
If there is an $n$-vertex weighted graph $G = (V,w)$ with $\cd(G) = k$
then $D_\lin(\MINCUT_n) \ge k$.
\end{corollary}

%






































































































































































































































































































































































































































































































































































































































































































































































































\setcounter{section}{6}\setcounter{theorem}{36}\section{$\ell_1$-approximate cut dimension}
\label{sec:l1-approx}
\label{sec:approx_construction}
In this section, we use the $\ell_1$-approximate cut dimension method to show \cref{thm:2n-2} that
for any $k \in \N$ and $n = 3k+1$, it holds that $D_\lin(\MINCUT_n) \ge 2n-2$.

Let $K_4$ be the complete graph on $4$ vertices with all edge weights equal to $1$.
The theorem will follow from showing that the $\ell_1$-approximate cut dimension of
the direct union of $k$ copies of $K_4$ has $\ell_1$-approximate cut dimension $6k$.
We start with the base case $k=1$ to build up the notation and intuition that will be needed
for the general case.  The following definition and fact will be useful.

\begin{definition}[Strictly diagonally dominant]
Let $A \in \R^{n \times n}$ be a matrix.  We say that the $\ith$ row of $A$ is \emph{strictly diagonally dominant}
if $|A(i,i)| > \sum_{j \ne i} |A(i,j)|$.  We say that $A$ is strictly diagonally dominant iff all of its rows are.
\end{definition}

It is well known that a strictly diagonally dominant matrix has full rank.
One way to prove this is via the following fact, which we will make use of in the proof of \cref{thm:2n-2}.
\begin{fact}
\label{fact:sdd}
Let $A \in \R^{n \times n}$ be a matrix whose $\ith$ row is strictly diagonally dominant.
If $A u = \zero$ for a vector $u \ne \zero$ then $|u_i| < \|u\|_\infty$.
\end{fact}

\begin{proof}
Suppose for a contradiction that for some $u \ne \zero$ it holds that $A u = \zero$ and $|u_i| = \|u\|_\infty$
where the $\ith$ row of $A$ is strictly diagonally dominant.
By normalizing and flipping the sign of $u$ if necessary we may assume $\|u\|_\infty =1$ and $A(i,i) u_i = |A(i,i)|$.  Thus
\begin{align*}
\sum_j A(i,j) u_j &= |A(i,i)| + \sum_{j \ne i} A(i,j) u_j \ge |A(i,i)| - \sum_{j \ne i} |A(i,j)| > 0 \enspace ,
\end{align*}
a contradiction.
\end{proof}


\subsection{$\ell_1$-approximate cut dimension of $K_4$}
\begin{figure}
\centering
\begin{tikzpicture}
\GraphInit[vstyle=Normal]
\tikzset{VertexStyle/.style = {draw, shape = circle, minimum size = 30pt}}
\Vertex[Math, x=0,y=0,L=v]{V}
\Vertex[Math,x=-2,y=2, L=a]{A}
\Vertex[Math,x=-4,y=0, L=b]{B}
\Vertex[Math,x=-2,y=-2, L=c]{C}
\Edge[label=3](V)(A)
\Edge[label=5, style={pos=.8}](V)(B)
\Edge[label=6](V)(C)
%
\Edge[label=1](A)(B)
\Edge[label=2,style={pos=.2}](A)(C)
\Edge[label=4](B)(C)
\end{tikzpicture}
\caption{The complete graph on $4$ vertices with all edge weights equal to $1$.  The labels on edges indicate the
ordering of edges used to represent cut vectors in the proof.}
\label{fig:K4}
\end{figure}

We label the vertices of $K_4$ by $a,b,c,v$, and
use the ordering of edges indicated in \cref{fig:K4}.
Let $X$ be the $7$-by-$6$ matrix whose rows correspond to the cut vectors of all the nontrivial cuts
\begin{equation}
\label{eq:baseX}
X =
\begin{bmatrix}
1 &  1 &  1 &  0 &  0 &  0 \\
1 &  0 &  0 &  1 &  1 &  0 \\
0 &  1 &  0 &  1 &  0 &  1 \\
0 &  0 &  1 &  0 &  1 &  1 \\
0 &  1 &  1 &  1 &  1 &  0 \\
1 &  0 &  1 &  1 &  0 &  1 \\
1 &  1 &  0 &  0 &  1 &  1 \\
\end{bmatrix}
\enspace .
\end{equation}
The cut vectors in $X$ are given in the order
\[
\Delta(\{a\}), \Delta(\{b\}), \Delta(\{c\}), \Delta(\{a,b,c\}), \Delta(\{a,b\}), \Delta(\{a,c\}), \Delta(\{b,c\}) \enspace.
\]
The first 4 rows correspond to star cuts which are minimum cuts of weight 3 in $K_4$.  The last three rows correspond to cuts which have weight 4 in $K_4$.
Thus to show a lower bound of $6$ on the number of linear queries needed to compute the minimum
cut of a 4 vertex graph, we need to show that the $w=(1,1,1,1,1,1), c=(0,0,0,0,1,1,1)$ one-sided $\ell_1$ approximate rank of $X$ is 6.

\begin{claim}
Let $w = \1 \in \R^6$, and $c =(0,0,0,0,1,1,1)$.  The $(w,c)$ one-sided $\ell_1$ approximate
rank of $X$ is $6$.
\end{claim}

\begin{proof}
The rank of $X$ at most $6$ as this is the number of columns, which takes care of the upper bound.

Now consider the lower bound.  To do this we need to lower bound the rank of the matrix
\[
Z=
X -
\begin{bmatrix}
\zero_{4,2} & \zero_{4,2} & \zero_{4,2} \\
A_1 & A_2 & A_3
\end{bmatrix}
\]
where each of $A_1, A_2,A_3 \ge 0$ are $3$-by-$2$ matrices and every row of $A_1 + A_2 + A_3$ sums to at most $1$.
As the first 4 rows of $X$ correspond to vectors of minimum cuts, no error is allowed on the first 4 rows.

The first 4 rows of $Z$ are equal to the first 4 rows of $X$, as there is no perturbation allowed on these rows.  By doing
elementary row operations on the first four rows, which do not change the rank, we can transform the first four rows of $Z$ into the reduced row echelon form
of $X(1:4, :)$.  Thus we arrive at the following matrix.
\[
\begin{bmatrix}
1 &  0 &  0 &  0 &  0 & -1 \\
0 &  1 &  0 &  0 & -1 &  0 \\
0 &  0 &  1 &  0 &  1 &  1 \\
0 &  0 &  0 &  1 &  1 &  1 \\
0 &  1 &  1 &  1 &  1 &  0 \\
1 &  0 &  1 &  1 &  0 &  1 \\
1 &  1 &  0 &  0 &  1 &  1 \\
\end{bmatrix}
-
\begin{bmatrix}
\zero_{4,2} & \zero_{4,2} & \zero_{4,2} \\
A_1 & A_2 & A_3
\end{bmatrix}
\enspace .
\]
Now we do column operations to zero out the entries in the first four rows and last two columns.  For a $m$-by-$2$ matrix $A$
we will use the notation $A^\circ$ to denote the matrix $A$ with the order of the columns swapped.
We arrive at
\[
\begin{bmatrix}
1 &  0 &  0 &  0 &  0 & 0 \\
0 &  1 &  0 &  0 & 0 &  0 \\
0 &  0 &  1 &  0 &  0 &  0 \\
0 &  0 &  0 &  1 &  0 &  0 \\
0 &  1 &  1 &  1 &  0 &  -2 \\
1 &  0 &  1 &  1 &  -2 &  0 \\
1 &  1 &  0 &  0 &  2 &  2 \\
\end{bmatrix}
-
\begin{bmatrix}
\zero_{4,2} & \zero_{4,2} & \zero_{4,2} \\
A_1 & A_2 & A_1^\circ -A_2 - A_2^\circ + A_3
\end{bmatrix}
\enspace .
\]
Finally, we can do row operations to zero out the first four columns in the last three rows.
\[
\begin{bmatrix}
1 &  0 &  0 &  0 &  0 & 0 \\
0 &  1 &  0 &  0 & 0 &  0 \\
0 &  0 &  1 &  0 &  0 &  0 \\
0 &  0 &  0 &  1 &  0 &  0 \\
0 &  0 &  0 &  0 &  0 &  -2 \\
0 &  0 &  0 &  0 &  -2 &  0 \\
0 &  0 &  0 &  0 &  2 &  2 \\
\end{bmatrix}
-
\begin{bmatrix}
\zero_{4,2} & \zero_{4,2} & \zero_{4,2} \\
\zero_{3,2} & \zero_{3,2} & A_1^\circ -A_2 - A_2^\circ + A_3
\end{bmatrix}
\enspace .
\]

The task has now reduced to showing the matrix
\[
Z'=
\begin{bmatrix}
2 & 0 \\
0 & 2 \\
-2 & -2
\end{bmatrix}
+ A_1 - A_2 - A_2^\circ + A_3^\circ
\]
has rank $2$ for any $A_1, A_2, A_3$ satisfying the constraints.  Let us simplify the matrix $A_1 - A_2 - A_2^\circ + A_3^\circ$.
First, let $A_1' = A_1 + A_3^\circ$.  Next, note that $D = A_2 + A_2^\circ$ has the property that $D(i,1) = D(i,2)$ for $i \in [3]$.
In the sequel we call this the \emph{partner property}.

As the row sum of $A_1' + A_2$ is at most $1$, unless $A_1'(1:2,1:2) = \zero_{2,2}$ and at least one row sum of $A_2(1:2,1:2)$ is equal to $1$ the
first two rows of $Z'$ will be strictly diagonally dominant.  If the first two rows of $Z'$ are strictly diagonally dominant then the
rank of $Z'$ must be $2$, thus we now handle the ``unless'' case.

First, suppose exactly one row sum of $A_2(1:2,1:2)$ is equal to $1$. Say without loss of generality it is the second one, thus
the first row of $Z'$ is strictly diagonally dominant.  Then for a sufficiently small $\varepsilon$ we can multiply the first column by $1-\varepsilon$
so that the first row remains strictly diagonally dominant and the second row becomes strictly diagonally dominant as well.  This
does not increase the rank and thus shows again that the rank of $Z'$ is $2$.

The remaining case is where both rows of $A_2(1:2,1:2)$ sum to one.  In this case by the partner property we have
\[
Z'(1:2,1:2) = \begin{bmatrix}
1&-1\\
-1&1
\end{bmatrix} .
\]
On the other hand, the last row of $Z'$ must have both entries $\le -1$.  Thus the determinant of the submatrix formed by the first row and the third is strictly negative and so $Z'$
has rank $2$.
\end{proof}


\subsection{Direct union of $K_4$ with itself}
\begin{figure}
\centering
\begin{minipage}{.55\textwidth}
\centering
\begin{tikzpicture}
\GraphInit[vstyle=Normal]
\tikzset{VertexStyle/.style = {draw, shape = circle, minimum size = 30pt}}
\Vertex[Math, x=0,y=0,L=v]{V}
\Vertex[Math,x=-2,y=2, L=a^{(1)}]{A}
\Vertex[Math,x=-4,y=0, L=b^{(1)}]{B}
\Vertex[Math,x=-2,y=-2, L=c^{(1)}]{C}
\Vertex[Math,x=2,y=2, L=a^{(2)}]{D}
\Vertex[Math,x=4,y=0, L=b^{(2)}]{E}
\Vertex[Math,x=2,y=-2, L=c^{(2)}]{F}
\Edge[label=3](V)(A)
\Edge[label=5, style={pos=.8}](V)(B)
\Edge[label=6](V)(C)
\Edge[label=9](V)(D)
\Edge[label=11, style={pos=.8}](V)(E)
\Edge[label=12](V)(F)
%
\Edge[label=1](A)(B)
\Edge[label=2,style={pos=.2}](A)(C)
\Edge[label=4](B)(C)
%
\Edge[label=7](D)(E)
\Edge[label=8,style={pos=.2}](D)(F)
\Edge[label=10](E)(F)
\end{tikzpicture}
\end{minipage}
\begin{minipage}{.4\textwidth}
\centering
\[
X^{(2)} =
\begin{bmatrix}
X & \zero_{7,6} \\
\zero_{7,6} & X
\end{bmatrix}
\]
\end{minipage}
\caption{Example of the direct union of two copies of $K_4$.  With the ordering of the edges given by the edge labels,
the matrix of cut vectors of the cuts $\Delta(\{a^{(i)}\}), \Delta(\{b^{(i)}\}), \Delta(\{c^{(i)}\}), \Delta(\{a^{(i)},b^{(i)},c^{(i)}\}), \Delta(\{a^{(i)},b^{(i)}\}), \Delta(\{a^{(i)},c^{(i)}\}), \Delta(\{b^{(i)},c^{(i)})\}$
for $i \in [2]$ becomes the matrix $X^{(2)}$ on the right.}
\label{fig:K42}
\end{figure}


Now we prove the general case.  The key to the proof is the following lemma.
\begin{lemma}
\label{lem:2k}
Let $k \in \N$ and $B$ be the $3k$-by-$2k$ matrix
\[
B = \begin{bmatrix} 2 \Id_{2k} \\ -2 \Id_k \otimes [1,1] \end{bmatrix} \enspace .
\]
For any matrices $3k$-by-$2k$ matrices $A_1, A_2$ satisfying the conditions
\begin{enumerate}
\item $A_1, A_2 \ge 0$
\item (partner property) For all $i \in [3k]$ and $j \in [k]$ it holds that $A_2(i,2j-1) = A_2(i,2j)$.
\item Every row of $A_1 + A_2/2$ sums to at most $1$
\end{enumerate}
it holds that $B +A_1 - A_2$ has rank $2k$.
\end{lemma}

\begin{proof}
The rank is at most $2k$ as that is the number of columns;  we focus on showing the columns are
linearly independent.

Let $Z = B + A_1 -A_2$.  We call the first $2k$ rows of $Z$ rows of type I, and the last $k$ rows
of type II.  If a type I row is not strictly diagonally dominant, we call it \emph{full}.  Notice that
a type I row $i$ is full if and only if the $\ith$ row of $A_1$ is zero and the $\ith$ row of $A_2$ sums to $2$.  In this
case, $Z(i,j) \le 0$ for every $j \ne i$ and it holds that $Z(i,i) = -\sum_{j \ne i} Z(i,j)$.
%
For $i \in [k]$ we call $2i-1$ and $2i$ \emph{partners}.

Suppose for contradiction there is a vector $\vecu \ne \zero$ such that $A\vecu = \zero$.  As $\vecu \ne \zero$ by
normalizing and multiplying by $-1$ as needed we may assume that $\|u\|_\infty = 1$ and $i$ is a coordinate with $\vecu(i)=1$.
By \cref{fact:sdd} the $\ith$ row of $Z$, which is a type I row, cannot be strictly diagonally dominant.  Thus the $\ith$ row
must be full.  Therefore for $Z(i,:) \vecu = 0$ to hold it must be the case that $\vecu(j) = 1$ for every $j$ where $A_2(i,j) > 0$.
Such a $j$ must exist as the $\ith$ row of $A_2$ sums to $2$.  So let $j$ be a coordinate with $A_2(i,j) > 0$ and let $j'$ be
the partner of $j$.  By the partner property we also have $A_2(i,j') > 0$ and therefore $\vecu(j) = \vecu(j') = 1$.

Now consider the type II row $\ell$ for which $B(\ell,j) = B(\ell,j') = -2$.
As $B(\ell,t) = 0$ for $t \not \in \{j,j'\}$ this means
\begin{align*}
Z(\ell,:) \vecu &= Z(\ell,j) + Z(\ell, j') + \sum_{t \not \in \{j,j'\}} Z(\ell,t) \vecu(t) \\
& \le B(\ell,j)+A_1(\ell,j) + B(\ell,j') + A_1(\ell,j') + \|\vecu\|_\infty \sum_{t \not \in \{j,j'\}} |Z(\ell,t)| \\
&\le -4 + \sum_{t} A_1(\ell,t) + A_2(\ell,t) \\
&\le -2 \enspace,
\end{align*}
and we have arrived at a contradiction.
 \end{proof}

With \cref{lem:2k} in hand we are now ready to prove \cref{thm:2n-2}.

\begin{proof}[Proof of \cref{thm:2n-2}]
Let $G^{(1)}, \ldots, G^{(k)}$ be $k$ copies of $K_4$ where the vertices in $G^{(i)}$ are labeled by $a^{(i)}, b^{(i)}, c^{(i)}, v^{(i)}$ for $i \in [k]$.
The graph $G$ is formed by taking the direct union of $G^{(1)}, \ldots, G^{(k)}$ at the vertices $v^{(1)}, \ldots, v^{(k)}$.  That is,
the vertices $v^{(1)}, \ldots, v^{(k)}$ are all identified by a common vertex denoted $v$.  See \cref{fig:K42} for an illustration of the graph
for $k=2$.

The cuts of $G$ we focus on are the $7k$ cuts given by
\[
\Delta(\{a^{(i)}\}), \Delta(\{b^{(i)}\}), \Delta(\{c^{(i)}\}), \Delta(\{a^{(i)},b^{(i)},c^{(i)}\}), \Delta(\{a^{(i)},b^{(i)}\}), \Delta(\{a^{(i)},c^{(i)}\}), \Delta(\{b^{(i)},c^{(i)}\} \enspace,
\]
for $i \in [k]$.  For any $i \in [k]$ the cuts $\Delta(\{a^{(i)}\}), \Delta(\{b^{(i)}\}), \Delta(\{c^{(i)}\}), \Delta(\{a^{(i)},b^{(i)},c^{(i)}\})$ achieve the minimum cut weight of $G$, which is 3,
and the cuts $\Delta(\{a^{(i)},b^{(i)}\}), \Delta(\{a^{(i)},c^{(i)}\}), \Delta(\{b^{(i)},c^{(i)}\}$ have weight 4.

With an ordering of the edges as exemplified in \cref{fig:K42}, the matrix of cut vectors of these cuts is $X^{(k)} = \Id_k \otimes X$,
where $X$ is the matrix from \cref{eq:baseX}.  In every nonzero block of $X^{(k)}$ the first four rows are minimum cuts with weight 3 and the last 3 rows are cuts with weight 4.
Let $c' = (0,0,0,0,1,1,1)$.  The theorem will follow from \cref{thm:approx_cut} by showing that the $w = \mathbf{1}_{6k}, c = \mathbf{1}_k \otimes c'$ one-sided $\ell_1$ approximate
rank of $X^{(k)}$ is $6k$.

To do this, we must show that $X^{(k)} - A$ has rank $6k$ for any matrix $A \ge 0$ which is all zero on any row of $I_k \otimes X$ corresponding
to a minimum cut, and where the row sum of $A$ is at most $1$ on any row of $I_k \otimes X$ corresponding to a cut of weight $4$.
In order to make reference to the base case, it will be useful to partition the columns into $k$ blocks of $6$ columns, where the $\ith$ block
is further partitioned into blocks of size $2$ represented by the $7k$-by-$2$ matrices $A_1^{(i)}, A_2^{(i)},A_3^{(i)}$.  In other words, we view
$A$ as follows
\[
A =
\begin{bmatrix}
A_1^{(1)} & A_2^{(1)} & A_3^{(1)} & \cdots & A_1^{(k)} & A_2^{(k)} & A_3^{(k)}
\end{bmatrix}
\enspace
\]
where each $A_j^{(i)}$ for $j \in [3], i \in [k]$ is a $7k$-by-2 matrix.

As in the base case, we begin by doing Gauss-Jordan elimination on the rows corresponding to mincuts of each $X$ block in $X^{(k)}$.  These operations
only touch rows corresponding to mincuts where $A$ is zero, thus they do not change $A$.  After these operations we arrive at the matrix $\Id_k \otimes X' - A$
where
\[
X' =
\begin{bmatrix}
1 &  0 &  0 &  0 &  0 & -1 \\
0 &  1 &  0 &  0 & -1 &  0 \\
0 &  0 &  1 &  0 &  1 &  1 \\
0 &  0 &  0 &  1 &  1 &  1 \\
0 &  1 &  1 &  1 &  1 &  0 \\
1 &  0 &  1 &  1 &  0 &  1 \\
1 &  1 &  0 &  0 &  1 &  1 \\
\end{bmatrix}
\]
Next, as in the base case, we do column operations to zero out the last two columns in the first four rows of each block of $X'$.  This
gives us the matrix $\Id_k \otimes X'' - A'$ where
\[
X'' =
\begin{bmatrix}
1 &  0 &  0 &  0 &  0 & 0 \\
0 &  1 &  0 &  0 & 0 &  0 \\
0 &  0 &  1 &  0 &  0 &  0 \\
0 &  0 &  0 &  1 &  0 &  0 \\
0 &  1 &  1 &  1 &  0 &  -2 \\
1 &  0 &  1 &  1 &  -2 &  0 \\
1 &  1 &  0 &  0 &  2 &  2 \\
\end{bmatrix}
\]
and the $\ith$ block of $A'$ looks like
\[
\begin{tabular}{c|c|c}
$[A_1^{(i)}$ & $A_2^{(i)}$ & $A_1^{(i)\circ} -A_2^{(i)} - A_2^{(i)\circ} + A_3^{(i)}] .$
\end{tabular}
\enspace.
\]
Here $A_1^{(i)\circ}$ denotes the matrix $A_1^{(i)}$ with the order of the columns swapped.
Finally, we use $X''(1:4,1:4)$ to zero out all other entries of $\Id_k \otimes X'' - A'$ in the first 4 columns of each block.
This brings us to the matrix $\Id_k \otimes X''' - A''$ where
\[
X''' =
\begin{bmatrix}
1 &  0 &  0 &  0 &  0 & 0 \\
0 &  1 &  0 &  0 & 0 &  0 \\
0 &  0 &  1 &  0 &  0 &  0 \\
0 &  0 &  0 &  1 &  0 &  0 \\
0 &  0 &  0 &  0 &  0 &  -2 \\
0 &  0 &  0 &  0 &  -2 &  0 \\
0 &  0 &  0 &  0 &  2 &  2 \\
\end{bmatrix}
\]
and the $\ith$ block of $A''$ is
\[
\begin{tabular}{c|c|c}
$[\zero_{7k,2}$ & $\zero_{7k,2}$ & $A_1^{(i)\circ} -A_2^{(i)} - A_2^{(i)\circ} + A_3^{(i)}]$
\end{tabular}
\enspace.
\]
Again, each of $A_1^{(i)}, A_2^{(i)}, A_3^{(i)}$ is zero on rows corresponding to minimum cuts.  Thus by multiplying the last two columns of
each block by $-1$ and permuting rows and columns we can transform $\Id_k \otimes X''' - A''$ into the form
\[
\begin{bmatrix}
\Id_{4k} & \zero_{4k,2k} \\
\zero_{3k,4k} & B +A_1 - A_2
\end{bmatrix}
\]
where $B,A_1, A_2$ satisfy the conditions of \cref{lem:2k}.  Thus a rank lower bound of $6k$ follows from the lower
bound of $2k$ on the rank of $B+A_1-A_2$ given in \cref{lem:2k}.
\end{proof}

%


























































































































































































































































































































































































































